% Generated draft: requires mathematical and source review.
\documentclass[11pt]{article}
\usepackage[margin=25mm]{geometry}
\usepackage{fontspec}
\setmainfont{Latin Modern Roman}
\usepackage{amsmath,amssymb,mathtools,unicode-math}
\setmathfont{Latin Modern Math}
\usepackage{enumitem,needspace,xurl,hyperref}
\setlength{\parindent}{0pt}
\setlength{\parskip}{5pt}
\setlength{\emergencystretch}{3em}
\newcommand{\N}{\mathbb{N}}
\newcommand{\Z}{\mathbb{Z}}
\newcommand{\Q}{\mathbb{Q}}
\newcommand{\R}{\mathbb{R}}
\newcommand{\C}{\mathbb{C}}
\newcommand{\F}{\mathbb{F}}
\newcommand{\A}{\mathbb{A}}
\newcommand{\PP}{\mathbb{P}}
\newcommand{\E}{\mathbb{E}}
\newcommand{\eps}{\varepsilon}
\newcommand{\dd}{\,\mathrm{d}}
\begin{document}
\section*{UnsolvedMath ID: 30006460}
\textbf{Source problem code:} catalogue entry
\subsection{Definitions and mathematical statement}
Let \(q\) denote an odd prime power, as specified in the supplied literature assessment. The intended setting concerns moments of function-field \(L\)-functions, with \(q\) held fixed in an asymptotic process. The supplied record does not identify the family of \(L\)-functions, the moment orders, the averaging measure, the asymptotic parameter, or the proposed formula and error term. Consequently, no precise moment conjecture or theorem can be recovered from this record.

\textbf{Unrecoverable mathematical statement.} The passage headed ``Conjecture 3'' is a discussion of methods and reported results, rather than the statement of that conjecture. According to the supplied source-verification note, the cited work proves a fixed-function-field moment formula for sufficiently large odd prime powers. The supplied dated assessment describes the full fixed-\(q\) statement for every odd prime power as unresolved. Neither the formula nor its hypotheses and quantifiers are supplied, so this description cannot be promoted to a standalone conjecture. In particular, the threshold for ``sufficiently large'' and its dependence on other parameters are unspecified.

The passage attributes Frobenius equidistribution to a theorem of Deligne and describes using large \(q\) to suppress contributions from higher-degree homology while computing lower-degree homology by homological stability. It claims that, for smooth \(X\), Frobenius eigenvalues on the object denoted \(H^i(X)\) have absolute values ``between \(q^{-i/2}\) and \(q^{-i}\).'' This is retained solely as a supplied claim: the base field, homology or cohomology theory, coefficients, Frobenius convention, normalization, and additional hypotheses are not specified. The passage alternates between cohomological notation and the word ``homology,'' preventing an unambiguous mathematical interpretation of this bound. Its discussion also contrasts fixing a particular \(q\) with first taking \(q\to\infty\), without specifying the other limit.

The source additionally reports formulas for shifted moments and asymptotics for ratios, the latter subject to a bound on distance from the critical line and attributed to Wang~[6]. The shifts, ratio parameters, distance condition, and asymptotic formulas are absent. These remarks are qualifications of the reported results, not independently specified mathematical statements.
\subsection{Short English statement}
This record preserves a discussion of a reported function-field moment result, but omits the conjecture and the moment formula themselves. Supplied curation reports results for sufficiently large odd prime powers and an unresolved extension to every fixed odd prime power. The averaging family, limits, hypotheses, and formulas must be recovered before a precise statement can be written. The draft requires human mathematical review.
\subsection{Sources}
\begin{itemize}
\item Local UnsolvedMath registry, permanent ID \(30006460\), title ``Reported Function-Field Moment Result,'' source problem code \texttt{catalogue entry}; supplied complete source record associated with \texttt{data/unsolvedmath/problems.json}. No proposer is identified; the supplied proposed year is 2025.
\item Supplied report citation: \emph{Analytic Number Theory (2026)}, Oberwolfach Reports, DOI 10.4171/owr/2025/51, pp.~2701--2764. The background describes Oberwolfach Reports as documenting workshops and their open-problem discussions. Supplied source and literature link, labeled ``Oberwolfach report 51/2025'': \url{https://doi.org/10.4171/owr/2025/51}.
\item Supplied literature reference, labeled ``Bergstr\"om--Diaconu--Petersen--Westerland'': \url{https://arxiv.org/abs/2302.07664}. No title or precise theorem reference is supplied.
\item Supplied literature reference, labeled ``Miller--Patzt--Petersen--Randal-Williams'': \url{https://arxiv.org/abs/2402.00354}. No title or precise theorem reference is supplied.
\item The passage attributes equidistribution and eigenvalue bounds to Deligne, refers to work numbered [2, 3], and attributes the ratios insight to Wang~[6]. The corresponding numbered bibliography is absent; no identification of those numbered references is inferred.
\item Literature assessment dated exactly \texttt{2026-08-22}: supplied dated curation, not independently verified. It labels the record ``partially solved'' and reports a restricted result using uniform twisted homological stability and sufficiently large \(q\), while describing the full fixed-\(q\) formulation for every odd prime power as unresolved. The background explicitly requests independent verification of time-sensitive status and cited claims.
\item The supplied statement-verification note says that the official report describes a proved formula for sufficiently large odd prime powers and extensions to shifted moments and ratios. It classifies the extracted passage as a proof discussion rather than a standalone open problem. The registry marks the statement ``unrecoverable,'' supplies no clean statement, and retains the original record for provenance. These are supplied verification claims; no source visit or independent verification is asserted here.
\end{itemize}
\end{document}
